\section{Quantitative evaluation: simulated displays}
\label{sec:results}

All numerical displays in this section use \emph{simulated data}, not measured
evaluator performance or collected human judgments. They instantiate the reporting
protocol (Appendix~\ref{app:simulation}); their assigned outcomes provide no
evidence for either research hypothesis.

\begin{table}[t]
\caption{\textbf{Simulated C1 reporting example; not experimental results.}
Each setting averages four equally weighted toy cohorts. Pair = pairwise
agreement; Top = top-region agreement; Bo4 = Best-of-4 accuracy (all in \%);
Reg = normalized rank regret (lower is better). Cohort assignments do not
assert actual label availability. Coverage and domain summaries accompany the
reproducible fixtures.}
\label{tab:results}
\centering\small
\setlength{\tabcolsep}{3.6pt}
\renewcommand{\arraystretch}{1.06}
% GENERATED SIMULATED DATA. Do not cite as experimental results.
\begin{tabularx}{\linewidth}{@{}P{3.8cm}*{8}{>{\raggedleft\arraybackslash}X}@{}}
\toprule
& \multicolumn{4}{c}{ID test} & \multicolumn{4}{c}{OOD test} \\
\cmidrule(lr){2-5}\cmidrule(l){6-9}
Evaluator & Pair & Top & Bo4 & Reg & Pair & Top & Bo4 & Reg \\
\midrule
\multicolumn{9}{@{}l}{\textcolor{pEvo}{\textbf{SIMULATED}}\quad Reused preferences} \\
Single metric & 66.5 & 67.5 & 44.8 & 0.31 & 64.1 & 65.2 & 43.5 & 0.32 \\
Metric ensemble & 69.9 & 70.9 & 51.2 & 0.26 & 66.8 & 68.0 & 45.3 & 0.30 \\
Static judge & 71.3 & 72.3 & 52.2 & 0.25 & 69.9 & 71.1 & 49.8 & 0.27 \\
Static tools & 74.0 & 75.2 & 57.1 & 0.23 & 71.3 & 72.9 & 52.8 & 0.24 \\
Prompt optimization & 74.1 & 75.5 & 56.8 & 0.22 & 71.9 & 72.9 & 54.4 & 0.24 \\
Program search & 74.8 & 76.1 & 57.7 & 0.21 & 73.0 & 74.5 & 55.5 & 0.23 \\
\rowcolor{pEvoWash}
IterEval & 75.3 & 76.9 & 61.2 & 0.20 & 72.7 & 74.2 & 55.3 & 0.23 \\
\addlinespace
\multicolumn{9}{@{}l}{\textcolor{pEvo}{\textbf{SIMULATED}}\quad New preferences} \\
Single metric & 66.2 & 67.0 & 45.1 & 0.30 & 63.3 & 64.3 & 40.5 & 0.34 \\
Metric ensemble & 68.3 & 69.6 & 47.8 & 0.29 & 65.9 & 66.9 & 43.8 & 0.32 \\
Static judge & 71.3 & 72.5 & 53.5 & 0.24 & 68.4 & 69.6 & 48.0 & 0.29 \\
Static tools & 72.4 & 73.6 & 54.9 & 0.24 & 70.0 & 71.2 & 51.3 & 0.26 \\
Prompt optimization & 74.3 & 75.5 & 56.8 & 0.22 & 70.6 & 71.8 & 51.3 & 0.26 \\
Program search & 74.3 & 75.7 & 56.5 & 0.22 & 71.6 & 72.9 & 53.4 & 0.24 \\
\rowcolor{pEvoWash}
IterEval & 75.8 & 77.1 & 59.9 & 0.20 & 72.2 & 73.6 & 53.4 & 0.24 \\
\bottomrule
\end{tabularx}

\end{table}

\paragraph{Alignment under distinct standards.}
Comparisons concern evolution within a task, not the effect of label origin.
Table~\ref{tab:results} separates the two construction
settings and the ID/OOD endpoints. Figure~\ref{fig:alignment}a--b retains
domain-level paired effects between ERA and the static tool-augmented seed on
jointly completed input groups. The rank-composition display in
Figure~\ref{fig:selection} distinguishes ordering agreement from selection.

\begin{figure}[t]
\centering
\input{figures/simulated_alignment}
\caption{\textbf{Simulated domain-level alignment; not experimental results.}
Panels (a) and (b) compare ERA with the static tool-augmented seed on ID and OOD
inputs. Diamonds are paired mean differences; whiskers are 95\% bootstrap
intervals over jointly completed synthetic input groups. Zero denotes no
difference. The two row groups indicate illustrative feedback-construction
settings, not permanent domain categories or measured data readiness.}
\label{fig:alignment}
\end{figure}

\paragraph{Acquisition as a separate test.}
Figure~\ref{fig:acquisition}a--b compares acquisition policies against random
sampling at equal human budgets. Held-out agreement after evolution is the
endpoint; each policy pays for the same reliability anchor. Annotation
granularity remains fixed and is varied separately in Figure~\ref{fig:ablation}b.

\begin{figure}[t]
\centering
\input{figures/simulated_acquisition}
\caption{\textbf{Simulated acquisition curves; not experimental results.}
Panels (a) and (b) compare five policies under reused and newly collected
preferences. Lines average 12 toy runs; bands mark the 10th--90th percentiles
for random and multi-source sampling, not confidence intervals. Budgets count
original-input annotation units, including a shared 10-unit random anchor.}
\label{fig:acquisition}
\end{figure}

\begin{samepage}
\paragraph{Feedback quality and unresolved judgments.}
Label origin does not determine reliability. Figure~\ref{fig:annotation-quality}a--c
separates valid-pair agreement, requested-rating abstention, and group consensus,
preserving their distinct denominators. Unreported reliability in reused
sources is marked unknown. C2 distributions and deployment-cost displays appear
in Figures~\ref{fig:refinement} and~\ref{fig:cost}; C3 remains entirely unmeasured.
\par
\end{samepage}
